\subsection*{Pooling the Linear Covariate Coefficients}

We estimate the treatment effect function $\alpha(\eta)$ via the semiparametric
partial linear model (PLM)
\begin{equation}\label{eq:plm}
  Y_i \;=\; X_i'\beta \;+\; h(\hat\eta_i) \;+\; D_i\,\alpha(\hat\eta_i) \;+\; \varepsilon_i,
\end{equation}
where $\hat\eta_i$ is the first-stage residual from regressing $Q_i$ on
$(1, X_i)$, $D_i = \mathbf{1}\{Q_i < 0\}$ is the treatment indicator,
$h(\cdot)$ is a B-spline approximation to the baseline conditional mean, and
$\alpha(\cdot)$ is a Gaussian kernel approximation to the treatment effect
function.
The key restriction in \eqref{eq:plm} is that $\beta$ is \emph{common} to
treated and control units---we call this the \emph{pooled} specification.
An alternative \emph{unpooled} specification adds $D_i X_i'\delta$ and allows
$\delta \neq 0$.

\paragraph{Why pooling is preferred.}
Both specifications are consistent when the true $\delta = 0$.
However, the unpooled specification is less efficient and can be substantially
biased in moderate samples for the following reason.
The treatment indicator $D_i = \mathbf{1}\{\gamma'X_i + \eta_i < 0\}$ induces
negative correlation between $X$ and $\hat\eta$ \emph{within the treated group}:
conditional on $D_i = 1$, a unit with large $X_i$ must have low $\eta_i$ to
fall below the cutoff.
This selection effect makes $D_i X_i$ correlated with $D_i \Phi_{\rm treat}(\hat\eta_i)$,
the columns that identify $\alpha(\eta)$.
Adding $D_i X_i$ to the design matrix therefore creates collinearity that
inflates the variance of $\hat\omega_{\rm treat}$, and in finite samples causes
the OLS to misallocate signal between $\hat\delta$ and $\hat\omega_{\rm treat}$.

\paragraph{Simulation evidence.}
We assess this via a Monte Carlo study using the DGP
$Q_i = \gamma X_i + \eta_i$, $Y_i = D_i W_i + \beta X_i + \nu_i$,
with $X_i \sim \mathcal{N}(0,1)$, $\eta_i$ bimodal, and
$\alpha(\eta) = 2e^\eta/(1+e^\eta)$ the true treatment effect.
Table~\ref{tab:pooling} reports the mean absolute error (MAE)
$\mathbb{E}|\hat\alpha(\eta) - \alpha(\eta)|$ averaged over the evaluation grid,
across $R = 200$ replications, for varying first-stage strength $\gamma$.

\begin{table}[h]
  \centering
  \caption{MAE of $\hat\alpha(\eta)$: pooled vs.\ unpooled ($\delta = 0$ by construction)}
  \label{tab:pooling}
  \small
  \begin{tabular}{lcccc}
    \hline
    & \multicolumn{2}{c}{$N = 5{,}000$} & \multicolumn{2}{c}{$N = 50{,}000$} \\
    \cmidrule(lr){2-3}\cmidrule(lr){4-5}
    $\gamma$ & Pooled & Unpooled & Pooled & Unpooled \\
    \hline
    0.5 & 0.188 & 0.190 & 0.206 & 0.206 \\
    1.0 & 0.114 & 0.118 & 0.125 & 0.127 \\
    1.5 & 0.077 & 0.083 & 0.073 & 0.075 \\
    2.0 & 0.065 & 0.069 & 0.048 & 0.050 \\
    2.5 & 0.059 & 0.063 & 0.036 & 0.037 \\
    \hline
  \end{tabular}
\end{table}

The pooled estimator dominates uniformly.
The advantage is modest---3--8\% lower MAE---and shrinks as $N$ grows,
consistent with both estimators being consistent.
The within-treated correlation $\mathrm{Corr}(X, \hat\eta \mid D=1)$ reaches
$-0.45$ at $\gamma = 1.5$ and the design-matrix collinearity
$\max_k |\mathrm{Corr}(DX, D\Phi_k(\hat\eta))|$ reaches $0.34$,
confirming the mechanism.
Note that the overall MAE level \emph{falls} with $\gamma$: a stronger first
stage identifies $\hat\eta$ more precisely, benefiting both estimators equally.

\paragraph{Implication for the GPA application.}
In the GPA data ($N \approx 42{,}500$, $p = 8$ covariates) the pooled and
unpooled $\hat\alpha$ differ by $0.084$ in the average treatment effect.
Under $\delta = 0$ this gap reflects cumulative collinearity across all eight
$D X_j$ columns, each individually correlated with $D\Phi_{\rm treat}$ at up
to $0.55$.
We therefore adopt the pooled specification \eqref{eq:plm} throughout.
